% Синтетический кейс: теоремы через ntheorem (без amsthm), в т.ч. \newframedtheorem.
\documentclass[12pt]{article}
\usepackage[T2A]{fontenc}
\usepackage[utf8]{inputenc}
\usepackage[russian]{babel}
\usepackage{amsmath,amssymb}
\usepackage{framed}
\usepackage[amsmath,thmmarks,framed]{ntheorem}

\theoremstyle{plain}
\newtheorem{theorem}{Теорема}[section]
\newtheorem{lemma}[theorem]{Лемма}
\theorembodyfont{\upshape}
\newtheorem{definition}[theorem]{Определение}
\newframedtheorem{framedclaim}[theorem]{Утверждение}
\theoremstyle{nonumberplain}
\theoremsymbol{\ensuremath{\square}}
\newtheorem{proof}{Доказательство}

\begin{document}

\section{Теоремы ntheorem}

\begin{definition}\label{def:n}
Определение ntheorem.
\end{definition}

\begin{theorem}[Главная]\label{thm:n}
Теорема ntheorem.
\end{theorem}

\begin{proof}
По определению~\ref{def:n}.
\end{proof}

\begin{lemma}
Лемма без метки.
\end{lemma}

\begin{framedclaim}
Утверждение в рамке.
\end{framedclaim}

\end{document}
